\section{Security-Aware Prompt Representation}

\model first turns a programming request into factors that can be named,
scoped, and edited while preserving the requested task. This makes an implicit
security requirement available as an intervention target and keeps task
behavior and surface presentation separately auditable.
Figure~\ref{fig:overview} previews this route from representation to hypothesis
formation and confirmation, with disjoint discover and held-out splits
separated by a pre-confirmation freeze.

\begin{figure}[t]
  \centering
  \resizebox{\linewidth}{!}{%
    \begin{tikzpicture}[
  font=\sffamily\footnotesize,
  input/.style={draw=black!65,rounded corners=6pt,align=center,
    minimum height=6.5mm,minimum width=47mm,fill=white},
  stage/.style={draw=black!80,rounded corners=2pt,align=center,
    text width=68mm,minimum height=12mm,fill=black!4,inner sep=2.2mm},
  record/.style={draw=black,rounded corners=2pt,align=center,
    text width=68mm,minimum height=11mm,fill=black!9,inner sep=2.2mm},
  downstream/.style={draw=black!60,densely dashed,rounded corners=2pt,
    align=center,text width=68mm,minimum height=9mm,fill=white,inner sep=2mm},
  flow/.style={-{Latex[length=2mm]},semithick},
  downstreamflow/.style={-{Latex[length=2mm]},semithick,densely dashed},
  note/.style={font=\sffamily\scriptsize,align=left,text=black!75}
]
  \node[input] (prompt) {Natural-Language Prompt};

  \node[stage,below=4mm of prompt] (repr) {
    \textbf{Stage I: Build Prompt TSG}\\[-0.2mm]
    typed context, requirement state, and valid operations\\[-0.2mm]
    \textit{Output: Typed Prompt Representation}};

  \node[stage,below=4mm of repr] (select) {
    \textbf{Stage II: Discover \& Freeze}\\[-0.2mm]
    TSG constraints, FCI, and semantic-cluster resampling\\[-0.2mm]
    \textit{Output: Frozen Requirement Hypothesis}};

  \node[stage,below=4mm of select] (evaluate) {
    \textbf{Stage III: Evaluate Effects}\\[-0.2mm]
    matched prompt policies and blocked randomization\\[-0.2mm]
    \textit{Output: Requirement-Effect Evidence}};

  \node[record,below=4mm of evaluate] (record) {
    \textbf{Requirement Evidence Record}\\[-0.2mm]
    context + operation + selection + effect + uncertainty + scope};

  \node[downstream,below=4mm of record] (expert) {
    \textbf{Expert Evaluation---RQ4}\\[-0.2mm]
    perceived quality and usefulness};

  \draw[flow] (prompt) -- (repr);
  \draw[flow] (repr) -- node[note,right=2mm] {derive templates} (select);
  \draw[flow] (select) -- node[note,right=2mm] {freeze before held-out outcomes} (evaluate);
  \draw[flow] (evaluate) -- (record);
  \draw[downstreamflow] (record) -- node[note,right=2mm] {separate study} (expert);
\end{tikzpicture}

  }
  \caption{\model converts prompt semantics into frozen intervention policies
  and evaluates their security effects on held-out tasks. Prompt TSG edges provide
  typed structure rather than causal conclusions; generated code enters only
  the independent Oracle and functional-evaluation path in the confirmation
  stage.}
  \Description{Three-stage \model workflow. The discover split contains
  security-aware prompt representation and stability-guided observational
  prioritization. A policy-freeze boundary separates it from randomized
  confirmation on held-out tasks using multi-realization prompt arms, block
  randomization, generated-code evaluation, and assigned-arm ITT estimation.}
  \label{fig:overview}
\end{figure}

\subsection{Typed Prompt Semantics}

For each prompt, \model separates three semantic families. Task-function
factors describe required operations and functional constraints; safety-control
factors describe guards, trust boundaries, assets, and security requirements;
and presentation-control factors describe wording and formatting used for
matched controls.

We encode these families in a bounded, typed task-security graph over the
prompt, which we call the \prompttsg.
Its relations encode prompt task/security structure---for example, that a
validation requirement is scoped to a sensitive input---not causal edges.
Every intervenable identity, scope, state, graph template, query, and permitted
ADD/REMOVE operation comes from an immutable versioned catalog. Extracted facts
must satisfy catalog closure and carry exact evidence spans into the prompt.
Deterministic validation checks those spans, family compatibility, and
duplicate or conflicting identities. Accepted records undergo canonical
node/edge identification, ordering, serialization, round-trip validation, and
digesting. The states \textsc{present}, \textsc{absent},
\textsc{not-applicable}, and \textsc{unresolved} remain distinct.

The intervention language distinguishes a non-editable context-query
specification $C_q$ from an actionable-feature specification $f$.
The former expresses guard-independent applicability, such as an untrusted
input reaching a shell operation; the latter names exactly one actionable
feature that an ADD or REMOVE policy may edit. A confirmable composite maps one
$C_q$ to exactly one actionable feature. Task operations, sources, sinks, APIs,
trust boundaries, and flow motifs remain context or effect-modifier
coordinates and are never silently edited to make a candidate actionable.

\subsection{Run-Locked Extraction}

\model supports three ways to build the same catalog-bound representation. An
LLM can emit catalog-bound semantic facts for deterministic graph construction
(\path{LLM_FACTS_V1}); it can instead emit a catalog-bound typed graph for
strict validation and canonicalization
(\path{LLM_DIRECT_GRAPH_V1}); or the reviewed finite matcher and templates
can construct the graph deterministically
(\path{DETERMINISTIC_CATALOG_V1}). Before processing any prompt, a run
selects exactly one backend and keeps it fixed throughout; per-prompt switching
and outcome-guided candidate selection are prohibited.

The extraction policy fixes provider, model, template, catalog, schema,
decoding, and request policy. The extractor treats the prompt as inert data and
is isolated from the target, arm, generated code, Oracle records, and outcomes.
It emits one semantic candidate. A transport failure may resend identical
request bytes, but the retry may not change semantic content or choose among
candidates for a favorable output.

\subsection{Local Causal Variables}

\model builds a local table for each configured CWE/security scope and model.
$W$ contains an approved minimal set of pre-treatment task metadata, $X^0$
contains non-deterministic natural-Prompt TSG queries, and $Y^0$ contains
independently committed discovery outcomes. Model identity remains an analysis
coordinate and provenance field. Raw text, generated code, code structure,
Oracle finding text, assigned arm $A$, and post-intervention $X^{A,R}$ are not
observational-table variables: generated code enters only the independent
Oracle and functional evaluator that produce $Y^0$.
